\section{Introduction}\label{sec:introduction}

Benchmark scores are the primary currency for comparing language models, yet training data contamination undermines the validity of these comparisons \citep{hendrycks2021mmlu,chang2025benchmarking}. The problem is not new, but its characterization remains superficial: existing work treats contamination as a binary property of models or benchmarks and proposes aggregate corrections that assume uniform effects across all test items. This framing masks critical structure. Within a single benchmark like MMLU, contamination may strongly inflate scores on mathematics items while leaving humanities items unaffected or producing negative effects. Aggregate corrections that average over such heterogeneity will systematically underestimate the problem in affected domains and overestimate it elsewhere. The result is neither accurate calibration nor actionable diagnostic information. Understanding contamination at domain granularity is therefore not merely desirable but necessary for any principled response to the problem.

Two developments make this gap urgent. First, \citet{sun2025emperor} systematically evaluate twenty contamination mitigation strategies (rephrasing, option shuffling, dynamic benchmark regeneration, and others) and demonstrate that none restores clean performance after contamination has occurred. This result forecloses the ``prevent contamination'' approach for models already trained on unrestricted web data. Second, the proliferation of open-weight models trained on trillion-token internet corpora means that contamination is now pervasive and cannot be precisely measured at the instance level \citep{oren2024proving}. Contamination-free benchmarks \citep{white2024livebench} offer a partial solution but sacrifice the comparability and institutional knowledge accumulated around established benchmarks. Dynamic regeneration of test items may also be contaminated within weeks as web crawlers index benchmark repositories \citep{dong_2024_cdd}. The field therefore needs post-hoc statistical tools that quantify contamination effects despite imperfect detection, and that provide domain-level diagnosis rather than aggregate point corrections.

Detection methods such as Min-K\%++ \citep{zhang_2024_minkpp}, ConStat \citep{dekoninck_2024_constat}, and Kernel Divergence Signatures \citep{choi_2025_kds} can flag contaminated items but stop at detection. They do not quantify how much scores are inflated, in which domains the inflation occurs, or whether effects are uniform across test items. Score adjustment methods like DCR \citep{xu2025dcr} provide point corrections without uncertainty quantification or domain-level granularity. IRT-based evaluation frameworks \citep{polo_2024_tinybenchmarks,habba_2025_growing_pains} apply psychometric ability estimation to improve evaluation efficiency and handle new models, but they ignore contamination entirely \citep{lalor_2024_irt_nlp}. These two research threads have developed in isolation: detection papers end at binary flags, while IRT papers assume all items are equally trustworthy. No existing method provides domain-level contamination \textit{diagnosis} with calibrated uncertainty. This gap prevents practitioners from understanding \textit{where} their benchmark scores are distorted, \textit{by how much}, and whether the distortion is uniform or heterogeneous across cognitive domains within a single benchmark.

We propose ContamCal, a two-stage diagnostic framework that bridges contamination detection and IRT-based evaluation. Stage~1 fuses multiple detection signals into continuous per-item exposure scores via training-free Reciprocal Rank Fusion \citep{cormack_2009_rrf}. This design avoids the meta-circularity of learned aggregators that would require the very contamination labels we lack. Stage~2 embeds these exposure scores as covariates in a contamination-augmented IRT model (2PL+$\gamma$) with per-domain contamination parameters that reveal where scores are inflated or deflated. The key empirical insight is a phenomenon we term \textit{dilution masking}: within MMLU, Math and Science domains show strong inflation ($\gamma \approx +0.97$ and $+0.84$ on the logit scale) while the pooled estimate across all eight domains is non-significant. The mechanism is straightforward: opposite-sign effects in non-STEM domains cancel the STEM inflation at the aggregate level. This cancellation is structurally analogous to DIF cancellation in educational measurement \citep{ercikan_2008_dif_cancellation}, and it renders any aggregate-level contamination analysis fundamentally unreliable when domain effects are heterogeneous.

Our contributions are as follows:
\begin{enumerate}
    \item \textbf{Domain-differential contamination characterization.} We demonstrate that MMLU contamination effects are heterogeneous across cognitive domains, with strong STEM inflation ($\gamma = +0.97$ for Math, $+0.84$ for Science) masked by null or negative non-STEM effects at the aggregate level. We term this phenomenon dilution masking and show it is a structural property of heterogeneous contamination, not a statistical artifact.
    \item \textbf{Detection signal failure modes.} We document a directional reversal in Min-K\%++ at model scales below 1.5B parameters (Cohen's $d = -1.55$ to $-1.98$), where the signal is actively misleading rather than merely ineffective. This finding has immediate practical consequences: practitioners applying Min-K\%++ to small models will reach the opposite of the correct conclusion.
    \item \textbf{Identifiability boundaries.} We derive Fisher Information bounds showing that reliable $\gamma$ recovery requires signal AUC $\geq 0.55$ and ensemble size $J \geq 50$. These thresholds reveal a structural mismatch between current LLM evaluation practice ($J = 8$--$30$) and the statistical requirements of contamination-augmented IRT.
    \item \textbf{Conditional calibration with value decomposition.} IRT ability estimation provides calibration benefit independent of $\gamma$; the contamination parameter serves a complementary diagnostic role that reveals which domains are distorted and by how much, rather than improving overall score accuracy. This decomposition clarifies the distinct purposes of calibration (how much to trust the score) and diagnosis (where and why the score is distorted).
\end{enumerate}